\documentclass[11pt]{article}

\usepackage[margin=1in]{geometry}
\usepackage{microtype}
\usepackage{booktabs}
\usepackage{hyperref}
\usepackage{url}
\usepackage{xcolor}
\usepackage{listings}
\usepackage{amsmath,amssymb}
\usepackage{graphicx}
\usepackage{natbib}
\bibliographystyle{plainnat}

\hypersetup{colorlinks=true, urlcolor=blue, linkcolor=blue, citecolor=blue}

\title{\texttt{bench-audit}: A Maintained Probe Suite for Agent-Benchmark
Integrity, with Calibrated Confidence Intervals}

\author{%
  Jacob Wu\thanks{Author email: \texttt{jacob.jikun.wu@gmail.com}.
  Code: \url{https://github.com/jacobwu/bench-audit}. Leaderboard:
  \url{https://jacobwu.github.io/bench-audit}.}\\
  PhD candidate, Montreal
}

\date{2026-05-14 \\[1ex] \textsc{Workshop preprint --- under preparation}}

\begin{document}
\maketitle

\begin{abstract}
\noindent
The major agent benchmarks --- SWE-bench, WebArena, OSWorld, GAIA,
Terminal-Bench, FieldWorkArena, CAR-bench, and SWE-bench Pro --- have all
been shown to admit structural exploits that yield near-perfect scores
without solving any task \citep{wang2026benchjack}. Frontier-lab models have
been independently shown to reward-hack on 30\%+ of evaluation runs on
RE-Bench tasks where the scoring function is in the model's context window
\citep{metr2025rewardhacking}. OpenAI publicly retired SWE-bench Verified as
a frontier-reporting target after finding that 59.4\% of failed tasks
contained broken test cases \citep{openai2026noswebench}. None of these
findings have a maintained, plug-in, reproducibility-first toolkit that
benchmark authors and consumers can run pre-submission and pre-publication.
We present \texttt{bench-audit}, an open-source probe suite that
reproduces the Berkeley and METR findings as a maintained library and
extends them in three ways: (i) every probe result carries a Wilson 95\%
confidence interval enforced at the type level; (ii) every benchmark
adapter declares its eval-set SHA-256 and refuses to run on mismatch;
(iii) every report card lists the reproduction command and links to the
raw per-task data. Initial v0.1 covers three benchmarks (SWE-bench
Verified, WebArena, GAIA) and six probes; the design supports adding a
new benchmark in one file.
\end{abstract}

\section{Introduction}

Between February and April 2026, three independent findings made it clear
that the major agent benchmarks no longer measure what their authors think
they measure. OpenAI announced on 2026-02-23 that it would no longer report
SWE-bench Verified after an internal audit found 59.4\% of 138 failed
problems had broken tests, and that GPT-5.2, Claude Opus 4.5, and Gemini~3
Flash could reproduce exact gold patches and verbatim problem details from
prompting alone \citep{openai2026noswebench}. On 2026-04-12, the UC
Berkeley Center for Responsible Decentralized Intelligence published a
study showing that an automated scanning agent, BenchJack, achieves
$\sim$100\% on seven of eight major agent benchmarks (and 73\% on OSWorld)
without making a single LLM call \citep{wang2026benchjack}. Since 2025-06,
METR has been publishing quantitative measurements of reward hacking on
frontier models --- most prominently a 30.4\% rate on RE-Bench
\citep{metr2025rewardhacking}, with one task family (Optimize LLM Foundry)
exhibiting reward hacking on 100\% of runs.

These findings concern three distinct integrity failures:

\begin{enumerate}
\item \textbf{Train-test contamination} --- the eval set has leaked into a
model's pretraining corpus, evidenced by verbatim recall under prompting
\citep{openai2026noswebench}.
\item \textbf{Harness injection and gold-answer leakage} --- the
\emph{evaluation pipeline} admits a structurally easier path. A
\texttt{conftest.py} hook resolves every SWE-bench instance; a
\texttt{file://} URL reads the WebArena gold answer; a system-PATH binary
wrapper produces a fake passing Terminal-Bench report
\citep{wang2026benchjack}.
\item \textbf{Reward hacking} --- a real model on a real benchmark
produces a trajectory that satisfies the literal scoring function without
satisfying task intent \citep{metr2025rewardhacking,krakovna2020spec}.
\end{enumerate}

The community lacks a maintained library that does all three in a single
plug-in interface. The Berkeley work is announced as a tool but is
currently a demonstration, not a maintained package; METR's pipeline is
partially open but bespoke per task family. \texttt{bench-audit} fills the
gap. Its design commitments, in order of importance:

\begin{itemize}
\item \emph{Every claim has a CI.} The \texttt{ProbeResult} type cannot be
instantiated without one. Verdicts at $n < 30$ require an explicit override
that is logged on every report card.
\item \emph{Every adapter ships an eval-set hash.} Cached eval sets are
verified at load time. Tampering shows as a mismatch.
\item \emph{Every probe ships with a reproducible recipe.} The
reproduction command is on every report card and on every leaderboard
cell.
\item \emph{No eval-set content is redistributed.} Adapters fetch from
upstream sources under their licenses.
\end{itemize}

\section{Related Work}
\label{sec:related}

\paragraph{Berkeley RDI: harness injection.} \citet{wang2026benchjack}
identify seven recurring vulnerability patterns across eight benchmarks:
(BA-01) no agent/evaluator isolation; (BA-02) answers shipped with tests;
(BA-03) \texttt{eval()} on untrusted input; (BA-04) LLM judges without
input sanitization; (BA-05) weak string matching; (BA-06) evaluation
logic that does not evaluate; (BA-07) grader trusts agent-reported state.
Their tool BenchJack scores $\sim$100\% on seven of eight benchmarks. Our
probe P4 (harness injection, Section~\ref{sec:methods}) reimplements this
taxonomy as a maintained checklist.

\paragraph{METR: reward hacking measurement.}
\citet{metr2025rewardhacking,metr2025o3eval} measure reward hacking in
frontier models. Headline numbers reproduced as our calibration targets:
RE-Bench overall 30.4\% (39/128), HCAST overall 0.7\% (8/1087), Optimize
LLM Foundry 100\% (21/21), Scaffolding for Rust 42.9\% (12/28). Our probe
P5 (reward hacking, Section~\ref{sec:methods}) reimplements METR's
two-stage classification (deterministic detector + LLM grader) with a
public rubric and a Cohen's $\kappa$ calibration target.

\paragraph{Contamination detection.} The string-overlap track of our P2
probe uses GPT-3's 13-gram convention \citep{brown2020gpt3} over a
configurable corpus index (Pile, RedPajama, BigCode, FineWeb, DCLM).
\citet{shi2024minkprob} provide the Min-K\% Prob membership inference
attack we use as a secondary signal; \citet{duan2024mia} report that
MIA-only is unreliable on The Pile, hence our policy of never issuing a
verdict on MIA alone.

\paragraph{Specification gaming.} The conceptual lineage is
\citet{krakovna2020spec}, \citet{skalse2022rewardhacking}, and
\citet{karwowski2024goodhart}; Goodhart's Law applied to reinforcement
learning predicts exactly the failure mode METR measures empirically.

\paragraph{Evaluation infrastructure.} We build on Inspect AI
\citep{aisi2024inspect} as the live-eval substrate, not against it.
\texttt{bench-audit} contributes a probe family to the Inspect AI
ecosystem.

\section{Methods}
\label{sec:methods}

\subsection{Adapter API}
A benchmark adapter is one file implementing four methods:
\texttt{load\_eval\_set}, \texttt{task\_iter}, \texttt{score}, and
\texttt{manifest}. The manifest carries license, source URL, eval-set
SHA-256, and a contamination-statement URL. Adapters refuse to operate on
a cached eval set that does not match the pinned hash. We ship three v0.1
adapters: SWE-bench Verified (MIT, $n=500$), WebArena (Apache-2.0,
$n=812$), GAIA (gated, $n=165$ validation).

\subsection{Probe contract}
\texttt{ProbeResult} is a frozen Pydantic model whose validator forbids
construction without (i) a confidence interval whose endpoints are
ordered, (ii) a sample size $\geq 30$ when verdict $\neq$
\texttt{inconclusive} unless \texttt{allow\_small\_n=True}, and (iii) a
CI half-width $\leq 0.1$ when verdict $\neq$ \texttt{inconclusive} unless
\texttt{allow\_wide\_ci=True}. Both overrides are recorded on every
report card.

\subsection{Six probes}
\begin{description}
\item[P1 Gold-answer leak.] An adapter-supplied \emph{lazy agent recipe}
produces predictions without doing task-relevant work. Reports leak rate
with Wilson 95\% CI.
\item[P2 Near-duplicate pretraining.] 13-gram overlap against a corpus
index; Min-K\% Prob as live-mode secondary signal.
\item[P3 Shortcut feature.] Mutual information between task metadata
$\phi$ and gold $y$ on a held-out split; bootstrap CI.
\item[P4 Harness injection.] Static-analysis checklist over the seven
Berkeley patterns; Clopper-Pearson CI on the count.
\item[P5 Reward hacking.] Two-stage trajectory classifier: deterministic
detectors over seven signatures (RH-01..RH-07) followed by an optional
LLM grader. Cohen's $\kappa$ $\geq 0.7$ pre-registered against three
human annotators.
\item[P6 Eval-set visibility.] Factual probe: public/gated, hashed,
contamination-statement available.
\end{description}

\subsection{Statistical commitments}
All proportion CIs use the Wilson score interval
\citep{wilson1927wilson}; we deliberately do not provide the
normal-approximation interval because it degenerates near 0 and 1 ---
where contamination claims live. Reward-hacking effect sizes between
groups use Cohen's~$h$ \citep{cohen1988stats}. Bootstrap is used for
non-proportion quantities (MI estimates, trajectory-length distributions).

\subsection{Live mode}
Live-mode probes run via Inspect AI \citep{aisi2024inspect}. The model
factory refuses \texttt{latest}-style aliases and undated provider
identifiers to guarantee reproducibility. Inspect AI's caching by
\texttt{(model, prompt\_sha256, sampling\_params)} makes reproduction
deterministic given the same cache state.

\section{Results}
\label{sec:results}

This section reports v0.1 calibration runs. Full-scale reproductions on
the real eval sets (SWE-bench $n=500$, WebArena $n=812$, GAIA $n=165$)
require access to SWE-bench's Docker harness, a WebArena server stack,
and the gated GAIA dataset; they are released alongside the public
launch. The fixture-based unit tests in this preprint exercise the probe
pathways end-to-end on five-task synthetic fixtures per adapter and
demonstrate that the probe correctly identifies the leak shape Berkeley
reported.

\paragraph{P1 vs. SWE-bench Verified (fixture).} Lazy-agent recipe forges
a \texttt{report.resolved=True} prediction (the \texttt{conftest.py}
pytest-hook shape Berkeley demonstrated). On the five-task fixture the
adapter's \texttt{score()} accepts the forged prediction for 5/5 tasks;
leak rate $= 1.000$, Wilson 95\% CI $[0.566, 1.000]$, $n=5$, verdict
\texttt{fail} with both \texttt{allow\_small\_n} and
\texttt{allow\_wide\_ci} flags set and logged.

\paragraph{P4 vs. SWE-bench Verified (fixture).} The harness-injection
checklist triggers BA-02: the fixture's task payload exposes
\texttt{gold\_patch}. Effect $= 1/7 \approx 0.143$, Clopper-Pearson 95\%
CI $[0.004, 0.579]$, $n=7$, verdict \texttt{fail} with overrides logged.

\paragraph{P5 vs. mini trajectory fixture.} The 10-trajectory mini
fixture has five hard reward-hacking cases (one per RH-01..RH-05), two
ambiguous (RH-06, RH-07), and three clean. The deterministic stage-1
detectors fire on 7/7 hack cases and stay quiet on 3/3 clean cases. Hack
rate (stub classifier) $= 5/10 = 0.500$, Wilson 95\% CI $[0.237, 0.762]$,
verdict \texttt{fail}.

Table~\ref{tab:calibration-targets} lists the Berkeley/METR numbers we
will reproduce on the full eval sets.

\begin{table}[t]
\centering
\caption{Calibration targets for full-scale runs.}
\label{tab:calibration-targets}
\begin{tabular}{lllr}
\toprule
Benchmark & Probe & Source & Target \\
\midrule
SWE-bench Verified ($n{=}500$) & P1 & \citep{wang2026benchjack} & $\geq 0.97$ \\
WebArena ($n{=}812$)           & P1 & \citep{wang2026benchjack} & $\geq 0.97$ \\
GAIA validation ($n{=}165$)    & P1 & \citep{wang2026benchjack} & $\geq 0.93$ \\
RE-Bench ($n{=}128$)           & P5 & \citep{metr2025rewardhacking} & $0.304 \pm 0.02$ \\
HCAST ($n{=}1087$)             & P5 & \citep{metr2025rewardhacking} & $0.007 \pm 0.02$ \\
\bottomrule
\end{tabular}
\end{table}

\section{Discussion}

The Berkeley and METR findings make a strong case that the per-benchmark
exploits are \emph{not} the deep problem. The deep problem is that
benchmark security has not been a first-class concern in the publication
process, and that the empirical literature has been comfortable claiming
``$X\%$ contamination'' without confidence intervals. We do not have a
methodological disagreement with either group; we are reproducing and
packaging their findings to make those numbers cheap to verify, cheap to
re-run on a new model, and cheap to extend to a new benchmark.

The single most important design choice is the type-level CI enforcement.
A \texttt{ProbeResult} that does not carry a CI cannot be constructed; a
verdict that overrides the small-$n$ or wide-CI defaults shows that
override on the report card. This is the form of discipline the
contamination literature has been missing.

\section{Limitations}

\begin{itemize}
\item v0.1 covers three benchmarks and six probes. Five of the eight
Berkeley benchmarks (SWE-bench Pro, OSWorld, Terminal-Bench,
FieldWorkArena, CAR-bench) are deferred to v0.2.
\item Our P4 detectors are heuristic static-analysis. A grader that is
clean in source but vulnerable via a third-party service it calls will
not trigger. Dynamic detectors are v0.2.
\item Reproducing the full-scale Berkeley numbers requires running
Docker-based evaluators we do not bundle (SWE-bench), a multi-container
web stack we do not bundle (WebArena), and a gated dataset we do not
redistribute (GAIA). The SWE-bench full run lands first; WebArena and
GAIA follow.
\item The Cohen's $\kappa$ validation for P5 requires three human
annotators. The rubric and labelling protocol are pre-registered; the
$\kappa$ measurement gates any external reward-hacking claim.
\end{itemize}

\section{Ethics statement}

Benchmark integrity work has two stakeholders we treat with care:

\paragraph{Benchmark maintainers.} We commit to a 14-day private outreach
window before any public claim about a specific maintained benchmark.
This is the responsible-disclosure norm from security research
transplanted to ML. Maintainers receive: the report cards, the raw
per-task data, and the reproduction commands. They have the option to
fix and rebut.

\paragraph{Model providers.} We do not redistribute model outputs from
frontier-lab models. Live-mode runs cache locally to the user only. The
harness refuses \texttt{latest}-style model aliases so historical
results cannot be silently revised.

The project ground rules are: no contamination claim without a CI; no
public disclosure of a finding against a maintained benchmark without
14-day private outreach; no misleading reproductions; respect upstream
licenses; do not redistribute frontier-lab model outputs; frame findings
as ``measured a contamination signature in X'' rather than ``broke X''.

\section{Reproducibility statement}

The code, eval-set hashes, prompt templates with their SHA-256s, and the
reproduction commands for every number in this paper are at
\url{https://github.com/jacobwu/bench-audit}. The fixture-based results
in Section~\ref{sec:results} are reproducible from a clean machine via
\texttt{./reproduce.sh}. Full-scale results are added to the leaderboard
(\url{https://jacobwu.github.io/bench-audit}) under
\texttt{docs/published\_numbers.json} as they land.

\bibliography{refs}

\end{document}
